\chapter{The \prg Runtime Interface}
\label{app:api}

This appendix is the working reference for the \prg framework calls used
throughout the book, together with the calling convention, the cartridge tool, and
the constants you will reach for most often. It catalogues the public interface
declared in \fname{components/prg32/include/prg32.h} and documented in the
platform's ABI notes \citep{prg32repo,prg32abi}. It is organised by task rather
than alphabetically, so that you can find a call by what you are trying to do.

\section{The calling convention}

Every framework function follows the standard \riscv calling convention, which is
what allows assembly games and C games to call the same functions
\citep{prg32abi}. Arguments are passed in \reg{a0}--\reg{a7} in order; the return
value comes back in \reg{a0}. Temporaries \reg{t0}--\reg{t6} may be destroyed by a
call; saved registers \reg{s0}--\reg{s11} are preserved across one. The return
address is in \reg{ra} and must be saved before a nested call; the stack pointer
\reg{sp} stays 16-byte aligned.

\begin{center}
\begin{tabularx}{\linewidth}{@{}llX@{}}
\toprule
Register & Role & Preserved across a call? \\
\midrule
\reg{a0}--\reg{a7} & arguments / return value & no \\
\reg{t0}--\reg{t6} & temporaries & no \\
\reg{s0}--\reg{s11} & saved values & yes \\
\reg{ra} & return address & no (save it yourself) \\
\reg{sp} & stack pointer & yes (16-byte aligned) \\
\bottomrule
\end{tabularx}
\captionof{table}{The \riscv register convention as \prg uses it \citep{prg32abi}.}
\end{center}

\section{Lifecycle and timing}

\begin{center}
\begin{tabularx}{\linewidth}{@{}lX@{}}
\toprule
Function & Purpose \\
\midrule
\code{prg32\_init()} & initialise the runtime; call once before the loop \\
\code{prg32\_ticks\_ms()} & milliseconds since start, for timing and animation \\
\code{prg32\_gfx\_present()} & display the frame you have drawn \\
\bottomrule
\end{tabularx}
\captionof{table}{Lifecycle and timing calls.}
\end{center}

\section{Input}

All buttons arrive as a single bitmask. Player two occupies the same layout shifted
into higher bits (\code{PRG32\_P2\_*}) \citep{prg32repo}.

\begin{center}
\begin{tabularx}{\linewidth}{@{}lX@{}}
\toprule
Function & Purpose \\
\midrule
\code{prg32\_input\_read()} & full player 1 / player 2 bitmask \\
\code{prg32\_input\_read\_player(p)} & one player's bits, normalised to the low bits \\
\code{prg32\_input\_read\_menu()} & both joysticks merged, for menu navigation \\
\code{prg32\_input\_wait\_released(mask)} & block until the selected bits are released \\
\bottomrule
\end{tabularx}
\captionof{table}{Input calls.}
\end{center}

\begin{center}
\begin{tabular}{@{}lr l lr@{}}
\toprule
Constant & Bit & & Constant & Bit \\
\midrule
\code{PRG32\_BTN\_LEFT}  & \code{0x01} & & \code{PRG32\_BTN\_A}      & \code{0x10} \\
\code{PRG32\_BTN\_RIGHT} & \code{0x02} & & \code{PRG32\_BTN\_B}      & \code{0x20} \\
\code{PRG32\_BTN\_UP}    & \code{0x04} & & \code{PRG32\_BTN\_SELECT} & \code{0x40} \\
\code{PRG32\_BTN\_DOWN}  & \code{0x08} & & (\code{PRG32\_BTN\_START} alias) & \code{0x40} \\
\bottomrule
\end{tabular}
\captionof{table}{Player-one button bits \citep{prg32repo}.}
\end{center}

\section{Graphics}

The viewport is 320 by 200 within a 320-by-240 display; colours are RGB565
\citep{prg32repo}.

\begin{center}
\begin{tabularx}{\linewidth}{@{}lX@{}}
\toprule
Function & Purpose \\
\midrule
\code{prg32\_gfx\_clear(color)} & fill the viewport with one colour \\
\code{prg32\_gfx\_pixel(x,y,color)} & set one pixel \\
\code{prg32\_gfx\_rect(x,y,w,h,color)} & filled rectangle \\
\code{prg32\_gfx\_text8(x,y,s,fg,bg)} & draw a text string \\
\code{prg32\_gfx\_present()} & present the finished frame \\
\code{prg32\_gfx\_lock()} / \code{\_unlock()} & enter / leave the graphics critical section \\
\bottomrule
\end{tabularx}
\captionof{table}{Core graphics calls.}
\end{center}

\begin{center}
\begin{tabular}{@{}lr l lr@{}}
\toprule
Colour & RGB565 & & Colour & RGB565 \\
\midrule
\code{PRG32\_COLOR\_BLACK} & \code{0x0000} & & \code{PRG32\_COLOR\_YELLOW}  & \code{0xffe0} \\
\code{PRG32\_COLOR\_WHITE} & \code{0xffff} & & \code{PRG32\_COLOR\_CYAN}    & \code{0x07ff} \\
\code{PRG32\_COLOR\_RED}   & \code{0xf800} & & \code{PRG32\_COLOR\_MAGENTA} & \code{0xf81f} \\
\code{PRG32\_COLOR\_GREEN} & \code{0x07e0} & & \code{PRG32\_COLOR\_BLUE}    & \code{0x001f} \\
\bottomrule
\end{tabular}
\captionof{table}{Named colour constants \citep{prg32repo}.}
\end{center}

\section{Sprites and collision}

\begin{center}
\begin{tabularx}{\linewidth}{@{}lX@{}}
\toprule
Function & Purpose \\
\midrule
\code{prg32\_sprite\_hitbox(ax,ay,aw,ah,bx,by,bw,bh)} & non-zero if two rectangles overlap \\
\code{prg32\_sprite\_anim\_frame(now,count,ms)} & which animation frame to show now \\
\code{prg32\_sprite\_draw\_frame(x,y,w,h,frames,frame,transparent)} & draw one frame of an animation \\
\code{prg32\_sprite\_draw\_8x8(x,y,bits,fg,bg)} & draw an 8-by-8 bit sprite \\
\bottomrule
\end{tabularx}
\captionof{table}{Sprite and collision calls \citep{prg32repo}.}
\end{center}

\section{Playfields and tiles}

The framework provides two scrolling tile layers for larger worlds
\citep{prg32repo}.

\begin{center}
\begin{tabularx}{\linewidth}{@{}lX@{}}
\toprule
Function & Purpose \\
\midrule
\code{prg32\_playfield\_put(layer,tx,ty,id)} & place a tile in the grid \\
\code{prg32\_playfield\_get(layer,tx,ty)} & read a tile from the grid \\
\code{prg32\_playfield\_scroll(layer,x,y)} & scroll a layer to a position \\
\code{prg32\_playfield\_parallax(layer,x\_q8,y\_q8)} & scroll with fixed-point parallax \\
\code{prg32\_playfield\_camera(x,y)} & set the shared camera position \\
\code{prg32\_playfield\_draw\_dual()} & draw both layers \\
\bottomrule
\end{tabularx}
\captionof{table}{Playfield and tile calls.}
\end{center}

\section{Platform-game physics}

The actor struct \code{prg32\_platform\_actor\_t} and its helpers implement
gravity, movement, and tile collision \citep{prg32repo}.

\begin{center}
\begin{tabularx}{\linewidth}{@{}lX@{}}
\toprule
Function & Purpose \\
\midrule
\code{prg32\_platform\_tile\_flags(id,flags)} & give a tile behaviour (solid, hazard, \dots) \\
\code{prg32\_platform\_actor\_init(\&a,layer,x,y,w,h)} & set up an actor \\
\code{prg32\_platform\_actor\_step(\&a,input,move,jump,grav,maxfall)} & one frame of physics; returns state flags \\
\code{prg32\_platform\_camera\_follow(\&a,dzx,dzy)} & scroll to keep the actor in view \\
\bottomrule
\end{tabularx}
\captionof{table}{Platform physics calls. Pass the actor by address (\code{\&}).}
\end{center}

\begin{center}
\begin{tabular}{@{}ll l ll@{}}
\toprule
Tile flag & Meaning & & Actor state & Meaning \\
\midrule
\code{...\_SOLID}    & blocks movement & & \code{...\_ON\_GROUND} & standing on ground \\
\code{...\_PLATFORM} & one-way platform & & \code{...\_HIT\_HEAD}  & bumped a ceiling \\
\code{...\_HAZARD}   & harmful tile & & \code{...\_HAZARD}    & touched a hazard \\
\code{...\_COLLECT}  & collectable & & \code{...\_COLLECT}   & collected an item \\
\bottomrule
\end{tabular}
\captionof{table}{Tile flags (\code{PRG32\_TILE\_FLAG\_*}) and actor state bits
(\code{PRG32\_PLATFORM\_*}) \citep{prg32repo}.}
\end{center}

\section{Console and sound}

\begin{center}
\begin{tabularx}{\linewidth}{@{}lX@{}}
\toprule
Function & Purpose \\
\midrule
\code{prg32\_console\_clear()} & clear the text console \\
\code{prg32\_console\_write(s)} & write a string to the console \\
\code{prg32\_console\_hex32(v)} & write a 32-bit value in hexadecimal \\
\code{prg32\_audio\_beep(hz,ms)} & play a tone (event sounds only) \\
\code{prg32\_audio\_note(midi,ms)} & play a MIDI note for a duration \\
\bottomrule
\end{tabularx}
\captionof{table}{Console and basic sound calls. Sound is audible on the board and
silent on QEMU.}
\end{center}

\section{Cartridges and scores}

\begin{center}
\begin{tabularx}{\linewidth}{@{}lX@{}}
\toprule
Function & Purpose \\
\midrule
\code{prg32\_cart\_default\_slot()} & saved default cartridge slot, or $-1$ \\
\code{prg32\_cart\_select\_default()} & load the saved default cartridge \\
\code{prg32\_score\_submit(game,player,score)} & record a local score \\
\code{prg32\_device\_demo\_run()} & run the built-in device demo \\
\bottomrule
\end{tabularx}
\captionof{table}{Selected cartridge and score calls \citep{prg32abi}.}
\end{center}

\section{The cartridge tool}

The \fname{tools/prg32\_game.py} script builds and deploys cartridges
\citep{prg32repo}. Its subcommands:

\begin{center}
\begin{tabularx}{\linewidth}{@{}lX@{}}
\toprule
Subcommand & Purpose \\
\midrule
\code{build} & compile a \code{.S} or \code{.c} source into a \code{.prg32} cartridge \\
\code{upload} & upload a cartridge to the board over HTTP (Wi-Fi) \\
\code{upload-qemu} & stage a cartridge into the QEMU flash image \\
\code{runtime} & print the runtime linker information \\
\code{doctor} & check local toolchain prerequisites \\
\bottomrule
\end{tabularx}
\captionof{table}{The cartridge tool's subcommands.}
\end{center}

The \code{build} subcommand's main options are \code{--out} (the output path,
required), \code{--name} (the cartridge name), \code{--entry-prefix} (the prefix of
your three functions), and \code{--firmware-elf} (the firmware ELF that resolves
the framework symbols). By default it targets the \riscv architecture
\code{rv32imc\_zicsr\_zifencei} with the \code{ilp32} ABI, the same 32-bit \riscv
profile your code is written for \citep{prg32repo}.

\section{The HTTP runtime API}

When the board's firmware is running and reachable, it exposes a small HTTP
interface useful for tooling and labs \citep{prg32repo}.

\begin{center}
\begin{tabularx}{\linewidth}{@{}lX@{}}
\toprule
Endpoint & Returns \\
\midrule
\code{GET /api/runtime} & firmware version, cartridge state, frame count, last input \\
\code{GET /api/games} & the installed cartridges \\
\code{POST /api/games?slot=cart0} & install a cartridge into a slot \\
\code{GET /api/screenshot.bmp} & the current 320-by-240 framebuffer as a BMP \\
\bottomrule
\end{tabularx}
\captionof{table}{Selected HTTP runtime endpoints \citep{prg32repo}.}
\end{center}
